\documentclass{beamer}
\usetheme{Madrid}
\usecolortheme{default}

% Define OSU orange color
\definecolor{osaorange}{RGB}{220,115,38}

% Set colors
\setbeamercolor{structure}{fg=osaorange}
\setbeamercolor{title}{fg=white,bg=osaorange}
\setbeamercolor{frametitle}{fg=white,bg=osaorange}

\usepackage{graphicx}
\usepackage{amsmath}
\usepackage{booktabs}
\usepackage{hyperref}
\usepackage{tikz}
\usetikzlibrary{shapes,arrows}

\title{H524 Introduction to Biostatistics}
\subtitle{Week 3: Theoretical Probability Distributions and Sampling Distributions}
\author{Dr. John Molitor}
\institute{}
\date{Fall 2025}

\begin{document}

\frame{\titlepage}

\begin{frame}
\frametitle{Outline}
\tableofcontents
\end{frame}

\section{Review from Week 2}

\begin{frame}
\frametitle{Quick Review: Probability Concepts}
\textbf{Last week we covered:}
\begin{itemize}
\item Basic probability rules and calculations
\item Conditional probability: $\Pr(A \mid B)$
\item Independence: $\Pr(A \cap B) = \Pr(A) \times \Pr(B)$
\item Diagnostic testing: sensitivity, specificity, PPV, NPV
\item Relative risk and odds ratios
\end{itemize}

\vspace{0.5cm}
\textbf{This Week:} Theoretical probability distributions and their role in statistical inference

\vspace{0.3cm}
\textbf{Reading:} Pagano \& Gauvreau Chapter 7 (7.1, 7.2, 7.4) and Chapter 8
\end{frame}

\section{Binomial Distribution (Ch 7.1)}

\begin{frame}
\frametitle{Why Probability Distributions?}
\textbf{From Week 2:} We calculated probabilities for specific events

\vspace{0.5cm}
\textbf{This Week:} Use theoretical models to describe entire families of outcomes

\vspace{0.5cm}
\textbf{Advantages of theoretical distributions:}
\begin{itemize}
\item Predict probabilities without collecting all possible data
\item Foundation for statistical inference (hypothesis testing, confidence intervals)
\item Standardized models applicable across many scenarios
\item Enable calculation of complex probabilities efficiently
\end{itemize}

\vspace{0.5cm}
\textbf{Two main types:}
\begin{itemize}
\item \textbf{Discrete distributions:} Countable outcomes (binomial, Poisson)
\item \textbf{Continuous distributions:} Uncountable outcomes (normal, t, chi-square)
\end{itemize}
\end{frame}

\begin{frame}
\frametitle{Binomial Distribution: Conditions}
\textbf{A binomial experiment must satisfy:}

\begin{enumerate}
\item \textbf{Fixed number of trials:} $n$ is predetermined
\item \textbf{Two possible outcomes:} Success or failure
\item \textbf{Constant probability:} $p$ remains the same for all trials
\item \textbf{Independent trials:} Outcome of one trial doesn't affect others
\item \textbf{Random variable:} $X$ = number of successes in $n$ trials
\end{enumerate}

\vspace{0.5cm}
\textbf{Notation:} $X \sim \text{Binomial}(n, p)$ or $X \sim B(n, p)$

\vspace{0.5cm}
\textbf{Medical Example:}
\begin{itemize}
\item $n = 20$ patients receive a vaccine
\item $p = 0.90$ probability vaccine protects (efficacy)
\item $X$ = number of patients protected
\end{itemize}
\end{frame}

\begin{frame}
\frametitle{Binomial Probability Mass Function}
\textbf{Probability of exactly $k$ successes:}

$$\Pr(X = k) = \binom{n}{k} p^k (1-p)^{n-k}$$

where $\binom{n}{k} = \frac{n!}{k!(n-k)!}$ is the binomial coefficient

\vspace{0.5cm}
\textbf{Parameters and Properties:}
\begin{itemize}
\item Expected value (mean): $\mu = E(X) = np$
\item Variance: $\sigma^2 = \text{Var}(X) = np(1-p)$
\item Standard deviation: $\sigma = \sqrt{np(1-p)}$
\item Range: $X \in \{0, 1, 2, \ldots, n\}$
\end{itemize}
\end{frame}

\begin{frame}
\frametitle{Binomial Example: Clinical Trial}
\textbf{Scenario:} A new treatment has 75\% success rate. We treat 10 patients.

\vspace{0.3cm}
\textbf{Question 1:} What's the probability exactly 8 patients respond?

\vspace{0.2cm}
$$\Pr(X = 8) = \binom{10}{8} (0.75)^8 (0.25)^2$$
$$= 45 \times 0.1001 \times 0.0625 = 0.282$$

\vspace{0.3cm}
\textbf{Question 2:} What's the probability at least 8 respond?

\vspace{0.2cm}
$$\Pr(X \geq 8) = \Pr(X=8) + \Pr(X=9) + \Pr(X=10)$$
$$= 0.282 + 0.188 + 0.056 = 0.526$$

\vspace{0.3cm}
\textbf{Expected number of successes:} $\mu = 10 \times 0.75 = 7.5$ patients
\end{frame}

\begin{frame}
\frametitle{Carcinogen Exposure Example}
\textbf{From Pagano \& Gauvreau:} 25\% of people exposed to a carcinogen become ill. Among 4 people with equal exposure:

\vspace{0.3cm}
\textbf{Calculate probabilities for 0, 1, 2, 3, 4 people becoming ill:}

\begin{center}
\begin{tabular}{cccc}
\toprule
$k$ & Calculation & Probability & Fraction \\
\midrule
0 & $\displaystyle\binom{4}{0}(0.25)^0(0.75)^4$ & 0.316 & 81/256 \\[0.4cm]
1 & $\displaystyle\binom{4}{1}(0.25)^1(0.75)^3$ & 0.422 & 108/256 \\[0.4cm]
2 & $\displaystyle\binom{4}{2}(0.25)^2(0.75)^2$ & 0.211 & 54/256 \\[0.4cm]
3 & $\displaystyle\binom{4}{3}(0.25)^3(0.75)^1$ & 0.047 & 12/256 \\[0.4cm]
4 & $\displaystyle\binom{4}{4}(0.25)^4(0.75)^0$ & 0.004 & 1/256 \\[0.4cm]
\midrule
 & Total & 1.000 & 256/256 \\
\bottomrule
\end{tabular}
\end{center}

\vspace{0.2cm}
\textbf{Probability at least one becomes ill:}\\
$1 - \Pr(X=0) = 1 - 0.316 = 0.684$
\end{frame}

\begin{frame}
\frametitle{Binomial Distribution in R}
\textbf{Four key functions:}

\vspace{0.3cm}
\begin{itemize}
\item \texttt{dbinom(k, n, p)} -- probability of exactly $k$ successes
\item \texttt{pbinom(k, n, p)} -- cumulative probability $\Pr(X \leq k)$
\item \texttt{qbinom(q, n, p)} -- quantile function (inverse of pbinom)
\item \texttt{rbinom(m, n, p)} -- generate $m$ random binomial values
\end{itemize}

\vspace{0.5cm}
\textbf{Example: Treatment with $n=10$, $p=0.75$}
\begin{itemize}
\item \texttt{dbinom(8, 10, 0.75)} returns 0.282
\item \texttt{pbinom(7, 10, 0.75)} returns $\Pr(X \leq 7) = 0.474$
\item \texttt{pbinom(7, 10, 0.75, lower.tail=FALSE)} returns $\Pr(X > 7) = 0.526$
\item \texttt{1 - pbinom(7, 10, 0.75)} also returns 0.526
\end{itemize}
\end{frame}

\section{Normal Distribution (Ch 7.2)}

\begin{frame}
\frametitle{From Discrete to Continuous}
\textbf{Continuous Random Variables:}
\begin{itemize}
\item Can take any value in an interval
\item Examples: height, weight, blood pressure, cholesterol
\item Probability described by a \textit{density curve}
\item $\Pr(a \leq X \leq b)$ = area under curve between $a$ and $b$
\end{itemize}

\vspace{0.5cm}
\textbf{Key properties of density curves:}
\begin{itemize}
\item Curve never goes below the x-axis
\item Total area under curve equals 1
\item $\Pr(X = c) = 0$ for any single point $c$
  \begin{itemize}
  \item Therefore $\Pr(a \leq X \leq b) = \Pr(a < X < b)$
  \end{itemize}
\end{itemize}
\end{frame}

\begin{frame}
\frametitle{The Normal Distribution}
\textbf{Most important continuous distribution in statistics}

\vspace{0.3cm}
\textbf{Characteristics:}
\begin{itemize}
\item Bell-shaped, symmetric curve
\item Completely determined by mean $\mu$ and standard deviation $\sigma$
\item Notation: $X \sim N(\mu, \sigma^2)$
\item Theoretical range: $-\infty$ to $+\infty$
\end{itemize}

\vspace{0.5cm}
\textbf{Important:} Not all bell-shaped curves are normal!
\begin{itemize}
\item t-distribution: similar shape but heavier tails
\item Chi-square: right-skewed for small df
\item Many biological variables are approximately normal
\end{itemize}
\end{frame}

\begin{frame}
\frametitle{The Empirical Rule (68-95-99.7)}
\textbf{For any normal distribution:}

\vspace{0.5cm}
\begin{itemize}
\item Approximately \textbf{68\%} of observations lie within \textbf{1 SD} of the mean
  $$\Pr(\mu - \sigma \leq X \leq \mu + \sigma) \approx 0.68$$

\item Approximately \textbf{95\%} lie within \textbf{2 SDs} of the mean
  $$\Pr(\mu - 2\sigma \leq X \leq \mu + 2\sigma) \approx 0.95$$

\item Approximately \textbf{99.7\%} lie within \textbf{3 SDs} of the mean
  $$\Pr(\mu - 3\sigma \leq X \leq \mu + 3\sigma) \approx 0.997$$
\end{itemize}

\vspace{0.3cm}
\textbf{Application:} Quick assessment of unusual observations
\begin{itemize}
\item Values beyond 2 SDs: unusual (only 5\% of data)
\item Values beyond 3 SDs: very unusual (only 0.3\% of data)
\end{itemize}
\end{frame}

\begin{frame}
\frametitle{Standard Normal Distribution}
\textbf{Special case:} $Z \sim N(0, 1)$ (mean = 0, SD = 1)

\vspace{0.3cm}
\textbf{Why is this useful?}
\begin{itemize}
\item We can standardize ANY normal distribution to $N(0,1)$
\item Only need one table (Table A.3) for all normal probabilities
\item Foundation for hypothesis testing and confidence intervals
\end{itemize}

\vspace{0.5cm}
\textbf{Standardization (Z-score transformation):}
$$Z = \frac{X - \mu}{\sigma}$$

\textbf{If $X \sim N(\mu, \sigma^2)$, then $Z \sim N(0, 1)$}

\vspace{0.3cm}
\textbf{Interpretation:} Z-score tells you how many SDs $X$ is from the mean
\begin{itemize}
\item $Z = 1.5$ means $X$ is 1.5 SDs above the mean
\item $Z = -2.0$ means $X$ is 2 SDs below the mean
\end{itemize}
\end{frame}

\begin{frame}
\frametitle{Working with the Standard Normal}
\textbf{Using Table A.3 (or R):}

\vspace{0.3cm}
If $Z \sim N(0, 1)$:
\begin{itemize}
\item $\Pr(Z \leq 0) = 0.5$ (symmetry)
\item $\Pr(Z < 1) = 0.8413$
\item $\Pr(Z < -1) = 0.1587$
\item $\Pr(Z > 1.96) = 1 - 0.975 = 0.025$
\item $\Pr(-1.96 < Z < 1.96) = 0.95$ (important for CIs!)
\end{itemize}

\vspace{0.5cm}
\textbf{R Functions:}
\begin{itemize}
\item \texttt{pnorm(z)} returns $\Pr(Z \leq z)$
\item \texttt{qnorm(p)} returns $z$ such that $\Pr(Z \leq z) = p$
\item \texttt{dnorm(z)} returns the density at $z$ (for plotting)
\end{itemize}
\end{frame}

\begin{frame}
\frametitle{Normal Distribution: Exam Scores Example}
\textbf{Biostatistics final exam scores:} $\mu = 70$, $\sigma = 10$

\vspace{0.3cm}
So $S \sim N(70, 10^2)$

\vspace{0.5cm}
\textbf{Question 1:} What proportion of students score above 85?

\vspace{0.2cm}
\textbf{Solution:}
\begin{enumerate}
\item Standardize: $Z = \frac{85 - 70}{10} = 1.5$
\item Find probability: $\Pr(S > 85) = \Pr(Z > 1.5)$
\item From table: $\Pr(Z < 1.5) = 0.9332$
\item Therefore: $\Pr(Z > 1.5) = 1 - 0.9332 = 0.0668$
\end{enumerate}

\vspace{0.3cm}
\textbf{Answer:} About 6.7\% of students score above 85

\vspace{0.3cm}
\textbf{R code:} \texttt{pnorm(85, mean=70, sd=10, lower.tail=FALSE)}
\end{frame}

\begin{frame}
\frametitle{Normal Distribution: Between Two Values}
\textbf{Same exam:} What proportion score between 60 and 85?

\vspace{0.3cm}
\textbf{Solution:}
\begin{enumerate}
\item Standardize lower bound: $Z_1 = \frac{60-70}{10} = -1.0$
\item Standardize upper bound: $Z_2 = \frac{85-70}{10} = 1.5$
\item $\Pr(60 < S < 85) = \Pr(-1 < Z < 1.5)$
\item $= \Pr(Z < 1.5) - \Pr(Z < -1.0)$
\item $= 0.9332 - 0.1587 = 0.7745$
\end{enumerate}

\vspace{0.3cm}
\textbf{Answer:} About 77.5\% score between 60 and 85

\vspace{0.3cm}
\textbf{R code:} \texttt{pnorm(85, 70, 10) - pnorm(60, 70, 10)}
\end{frame}

\begin{frame}
\frametitle{Finding Percentiles}
\textbf{Question:} What score is at the 90th percentile? (Top 10\%)

\vspace{0.2cm}
\textbf{Solution - Working backwards:}
\begin{enumerate}
\item We want $S_{90}$ such that $\Pr(S < S_{90}) = 0.90$
\item Standardize: $\Pr\left[Z < \frac{S_{90} - 70}{10}\right] = 0.90$
\item From table: $\Pr(Z < 1.28) = 0.90$
\item Therefore: $\frac{S_{90} - 70}{10} = 1.28$
\item Solve: $S_{90} = 70 + 1.28(10) = 82.8$
\end{enumerate}

\vspace{0.2cm}
\textbf{Answer:} Score of 82.8 is the 90th percentile

\vspace{0.2cm}
\textbf{R code:} \texttt{qnorm(0.90, mean=70, sd=10)} returns 82.8
\end{frame}

\begin{frame}
\frametitle{Medical Example: Cholesterol Levels}
\textbf{Adult cholesterol:} $\mu = 200$ mg/dL, $\sigma = 40$ mg/dL

\vspace{0.5cm}
\textbf{Clinical thresholds:}
\begin{itemize}
\item $< 200$ mg/dL: Desirable
\item $200-239$ mg/dL: Borderline high
\item $\geq 240$ mg/dL: High (treatment recommended)
\end{itemize}

\vspace{0.5cm}
\textbf{Question:} What proportion need treatment (cholesterol $\geq 240$)?

\vspace{0.3cm}
\textbf{Solution:}
\begin{itemize}
\item $Z = \frac{240 - 200}{40} = 1.0$
\item $\Pr(X \geq 240) = \Pr(Z > 1.0) = 1 - 0.8413 = 0.1587$
\end{itemize}

\vspace{0.3cm}
\textbf{Answer:} About 16\% of adults have high cholesterol requiring treatment
\end{frame}

\section{Other Important Distributions (Ch 7.4)}

\begin{frame}
\frametitle{Student's t-Distribution}
\textbf{Named after ``Student'' (William Sealy Gosset)}
\begin{itemize}
\item Brewer and statistician at Guinness
\item Published under pseudonym ``Student'' (1908)
\item Developed for small sample inference
\end{itemize}

\vspace{0.5cm}
\textbf{Properties:}
\begin{itemize}
\item Bell-shaped like normal, but heavier tails
\item Determined by degrees of freedom (df = n - 1)
\item As df increases, approaches standard normal
\item Used when $\sigma$ is unknown (estimated by $s$)
\end{itemize}

\vspace{0.5cm}
\textbf{When to use:}
$$T = \frac{\bar{X} - \mu}{s/\sqrt{n}} \sim t_{n-1}$$

\textbf{We'll use this extensively for confidence intervals and hypothesis tests!}
\end{frame}

\begin{frame}
\frametitle{Comparing Normal and t-Distributions}
\textbf{Key differences:}

\vspace{0.3cm}
\begin{center}
\begin{tabular}{lll}
\toprule
Feature & Normal $N(0,1)$ & t-distribution \\
\midrule
Shape & Bell-shaped & Bell-shaped \\
Tails & Thinner & Heavier \\
Parameters & None (standard) & Degrees of freedom \\
Use & $\sigma$ known & $\sigma$ unknown \\
Critical value (95\%) & 1.96 & Depends on df \\
\bottomrule
\end{tabular}
\end{center}

\vspace{0.5cm}
\textbf{Example critical values for 95\% confidence:}
\begin{itemize}
\item $Z_{0.975} = 1.96$ (always)
\item $t_{0.975, df=5} = 2.571$ (small sample)
\item $t_{0.975, df=30} = 2.042$ (moderate sample)
\item $t_{0.975, df=\infty} = 1.96$ (same as Z)
\end{itemize}
\end{frame}

\begin{frame}
\frametitle{Other Distributions (Brief Preview)}
\textbf{Chi-Square ($\chi^2$) Distribution:}
\begin{itemize}
\item Right-skewed, non-negative values
\item Used for: goodness-of-fit tests, contingency tables, variance testing
\item Parameters: degrees of freedom
\end{itemize}

\vspace{0.5cm}
\textbf{F-Distribution:}
\begin{itemize}
\item Ratio of two chi-square distributions
\item Used for: ANOVA, comparing variances
\item Parameters: two degrees of freedom (numerator and denominator)
\end{itemize}

\vspace{0.5cm}
\textbf{Poisson Distribution:}
\begin{itemize}
\item Discrete, for rare events
\item Used for: disease counts, rare adverse events
\item Parameter: rate $\lambda$ (mean = variance)
\end{itemize}

\vspace{0.3cm}
\textit{We'll encounter these in later weeks!}
\end{frame}

\section{Sampling Distributions (Ch 8)}

\begin{frame}
\frametitle{From Populations to Samples}
\textbf{Fundamental Statistical Problem:}
\begin{itemize}
\item \textbf{Population:} All subjects of interest (usually unknown)
  \begin{itemize}
  \item Has parameters: $\mu$ (mean), $\sigma$ (SD)
  \end{itemize}
\item \textbf{Sample:} Subset we actually observe
  \begin{itemize}
  \item Has statistics: $\bar{x}$ (sample mean), $s$ (sample SD)
  \end{itemize}
\end{itemize}

\vspace{0.5cm}
\textbf{Statistical Inference:}
\begin{itemize}
\item Use sample statistics to learn about population parameters
\item Account for sampling variability (uncertainty)
\item Quantify our confidence in conclusions
\end{itemize}

\vspace{0.5cm}
\textbf{Key Insight:} Different samples give different results. How much variation should we expect?
\end{frame}

\begin{frame}
\frametitle{Thought Experiment: OSU Student Heights}
\textbf{Imagine:}
\begin{itemize}
\item I take a sample of $n=100$ OSU students, measure heights
\item Compute sample mean: $\bar{x}_1 = 66.2$ inches
\item \textbf{You} take a different sample of $n=100$ students
\item Compute sample mean: $\bar{x}_2 = 65.8$ inches
\item \textbf{Everyone in class} takes a sample of $n=100$
\end{itemize}

\vspace{0.5cm}
\textbf{Questions:}
\begin{enumerate}
\item Will all sample means be the same? \textbf{No!}
\item Will they be close to each other? \textbf{Probably}
\item What's the distribution of all possible $\bar{x}$ values?
\end{enumerate}

\vspace{0.3cm}
\textbf{Answer:} The distribution of $\bar{x}$ is called the \textit{sampling distribution of the mean}
\end{frame}

\begin{frame}
\frametitle{Sampling Distribution of the Mean}
\textbf{Definition:} The probability distribution of all possible sample means from samples of size $n$

\vspace{0.5cm}
\textbf{Key Properties (for samples from ANY population):}
\begin{enumerate}
\item \textbf{Center:} $E(\bar{X}) = \mu$ (unbiased estimator)
  \begin{itemize}
  \item Sample mean centers on population mean
  \end{itemize}

\item \textbf{Spread:} $\text{SD}(\bar{X}) = \frac{\sigma}{\sqrt{n}}$ (standard error)
  \begin{itemize}
  \item Variability decreases as sample size increases
  \end{itemize}

\item \textbf{Shape:} Approximately normal for large $n$ (Central Limit Theorem)
  \begin{itemize}
  \item Even if population is NOT normal!
  \end{itemize}
\end{enumerate}

\vspace{0.5cm}
\textbf{Notation:} $\bar{X} \sim N\left(\mu, \frac{\sigma^2}{n}\right)$ for large $n$
\end{frame}

\begin{frame}
\frametitle{Standard Error vs Standard Deviation}
\textbf{CRITICAL DISTINCTION:}

\vspace{0.3cm}
\textbf{Standard Deviation ($\sigma$):}
\begin{itemize}
\item Variability of \textit{individual observations}
\item Describes spread in the population
\item Does NOT decrease with sample size
\end{itemize}

\vspace{0.3cm}
\textbf{Standard Error ($\text{SE} = \sigma/\sqrt{n}$):}
\begin{itemize}
\item Variability of \textit{sample means}
\item Describes precision of our estimate
\item Decreases as sample size increases
\item Smaller SE = more precise estimate
\end{itemize}

\vspace{0.3cm}
\textbf{Example:} Blood pressure with $\sigma = 20$ mmHg
\begin{itemize}
\item Individual SD: $\sigma = 20$ mmHg (always)
\item SE for $n=25$: $20/\sqrt{25} = 4$ mmHg
\item SE for $n=100$: $20/\sqrt{100} = 2$ mmHg
\end{itemize}
\end{frame}

\begin{frame}
\frametitle{Central Limit Theorem (CLT)}
\textbf{Most Important Theorem in Statistics!}

\vspace{0.5cm}
\textbf{Statement:}
For a random sample of size $n$ from ANY population with mean $\mu$ and finite variance $\sigma^2$:

$$\bar{X} \xrightarrow{d} N\left(\mu, \frac{\sigma^2}{n}\right) \text{ as } n \to \infty$$

Or equivalently:
$$Z = \frac{\bar{X} - \mu}{\sigma/\sqrt{n}} \xrightarrow{d} N(0, 1)$$

\vspace{0.5cm}
\textbf{In plain English:}
\begin{itemize}
\item Sample means are approximately normally distributed
\item Works for \textbf{ANY} population distribution (skewed, bimodal, etc.)
\item Approximation improves as $n$ increases
\item Rule of thumb: $n \geq 30$ often sufficient
\end{itemize}
\end{frame}

\begin{frame}
\frametitle{Why CLT Matters}
\textbf{The CLT enables statistical inference!}

\vspace{0.5cm}
\textbf{Applications:}
\begin{enumerate}
\item \textbf{Confidence Intervals:}
  \begin{itemize}
  \item We can use normal distribution to quantify uncertainty
  \item Even when population is not normal
  \end{itemize}

\item \textbf{Hypothesis Testing:}
  \begin{itemize}
  \item Calculate p-values using normal probabilities
  \item Test claims about population means
  \end{itemize}

\item \textbf{Sample Size Determination:}
  \begin{itemize}
  \item Predict precision of estimates before collecting data
  \item Plan studies efficiently
  \end{itemize}
\end{enumerate}

\vspace{0.5cm}
\textbf{Key Insight:} We don't need to know the population distribution to make inferences about the mean!
\end{frame}

\begin{frame}
\frametitle{CLT Example: IQ Scores}
\textbf{Population:} Los Angeles residents, IQ $\sim N(100, 15^2)$

Actually normal, but CLT still applies!

\vspace{0.3cm}
\textbf{Sample:} $n = 90$ residents, known variance $\sigma^2 = 100$ (so $\sigma = 10$)

\vspace{0.3cm}
\textbf{Observed:} Sample mean $\bar{x} = 105$

\vspace{0.5cm}
\textbf{Sampling distribution of $\bar{X}$:}
\begin{itemize}
\item Mean: $E(\bar{X}) = \mu = 100$
\item Standard error: $\text{SE} = \frac{10}{\sqrt{90}} = 1.054$
\item Distribution: $\bar{X} \sim N(100, 1.054^2)$
\end{itemize}

\vspace{0.3cm}
\textbf{Question:} How unusual is $\bar{x} = 105$?
\begin{itemize}
\item $Z = \frac{105 - 100}{1.054} = 4.74$ (!!)
\item $\Pr(Z > 4.74) < 0.00001$
\item Very strong evidence that LA mean IQ $> 100$
\end{itemize}
\end{frame}

\section{Confidence Intervals (Introduction)}

\begin{frame}
\frametitle{From Point Estimates to Intervals}
\textbf{Problem with point estimates:}
\begin{itemize}
\item Sample mean $\bar{x} = 105$ is our best guess for $\mu$
\item But we know $\bar{x}$ has sampling variability
\item How confident are we? What's the range of plausible values?
\end{itemize}

\vspace{0.5cm}
\textbf{Solution: Confidence Intervals}
\begin{itemize}
\item Provide a range of plausible values for $\mu$
\item Quantify our uncertainty
\item Account for sampling variability
\end{itemize}

\vspace{0.5cm}
\textbf{General Form:}
$$\text{Estimate} \pm \text{Margin of Error}$$
$$\bar{x} \pm (\text{critical value}) \times \text{SE}$$
\end{frame}

\begin{frame}
\frametitle{95\% Confidence Interval ($\sigma$ known)}
\textbf{When population SD ($\sigma$) is known:}

$$95\% \text{ CI for } \mu: \quad \bar{x} \pm 1.96 \times \frac{\sigma}{\sqrt{n}}$$

\vspace{0.5cm}
\textbf{Why 1.96?}
\begin{itemize}
\item For $Z \sim N(0,1)$: $\Pr(-1.96 < Z < 1.96) = 0.95$
\item Middle 95\% of standard normal distribution
\item Comes from $\bar{X} \sim N(\mu, \sigma^2/n)$ by CLT
\end{itemize}

\vspace{0.5cm}
\textbf{IQ Example:} $\bar{x} = 105$, $\sigma = 10$, $n = 90$
$$\text{95\% CI} = 105 \pm 1.96 \times \frac{10}{\sqrt{90}}$$
$$= 105 \pm 2.07 = (102.93, 107.07)$$

\textbf{Interpretation:} We are 95\% confident the true mean LA IQ is between 102.93 and 107.07
\end{frame}

\begin{frame}
\frametitle{95\% Confidence Interval ($\sigma$ unknown)}
\textbf{Reality:} We usually don't know $\sigma$!

\vspace{0.3cm}
\textbf{Solution:} Use sample SD ($s$) and t-distribution

$$95\% \text{ CI for } \mu: \quad \bar{x} \pm t_{0.975, n-1} \times \frac{s}{\sqrt{n}}$$

\vspace{0.5cm}
\textbf{Key differences:}
\begin{itemize}
\item Use $s$ instead of $\sigma$ (estimate from data)
\item Use $t_{n-1}$ critical value instead of 1.96
\item Accounts for extra uncertainty from estimating $\sigma$
\item \textbf{Assumes:} Population is approximately normal
\end{itemize}

\vspace{0.5cm}
\textbf{Example critical values:}
\begin{itemize}
\item $n = 10$: $t_{0.975, 9} = 2.262$ (wider than 1.96)
\item $n = 30$: $t_{0.975, 29} = 2.045$
\item $n = 100$: $t_{0.975, 99} = 1.984$ (close to 1.96)
\end{itemize}
\end{frame}

\begin{frame}
\frametitle{Tumor Growth Example}
\textbf{Clinical Trial:} 20 patients, 10 treatment, 10 control

\vspace{0.3cm}
\textbf{Control group tumor growth (mm):}\\
7, 10, 9, 8, 7, 6, 8, 9, 12, 13

\vspace{0.3cm}
\textbf{Treatment group tumor growth (mm):}\\
4, 6, 10, 8, 5, 3, 10, 8, 8, 10

\vspace{0.5cm}
\textbf{Control group:}
\begin{itemize}
\item $\bar{x}_C = 8.9$ mm, $s_C = 2.234$ mm, $n_C = 10$
\item $t_{0.975, 9} = 2.262$
\item 95\% CI: $8.9 \pm 2.262 \times \frac{2.234}{\sqrt{10}} = 8.9 \pm 1.598$
\item \textbf{95\% CI: (7.30, 10.50) mm}
\end{itemize}

\vspace{0.3cm}
\textbf{Treatment group:}
\begin{itemize}
\item $\bar{x}_T = 7.2$ mm, $s_T = 2.573$ mm
\item 95\% CI: \textbf{(5.36, 9.04) mm}
\end{itemize}
\end{frame}

\begin{frame}
\frametitle{Interpreting Confidence Intervals}
\textbf{Correct interpretation:}
\begin{itemize}
\item ``We are 95\% confident that the true population mean is between 7.30 and 10.50 mm''
\item The interval (7.30, 10.50) either contains $\mu$ or it doesn't
\item If we repeated this procedure many times, 95\% of intervals would contain $\mu$
\end{itemize}

\vspace{0.5cm}
\textbf{Common misinterpretations (WRONG):}
\begin{itemize}
\item $\times$ ``There's a 95\% probability $\mu$ is in (7.30, 10.50)''
  \begin{itemize}
  \item $\mu$ is fixed, not random; it's either in or out
  \end{itemize}
\item $\times$ ``95\% of the data is between 7.30 and 10.50''
  \begin{itemize}
  \item This is about the population mean, not individual observations
  \end{itemize}
\item $\times$ ``The sample mean has 95\% probability of being in (7.30, 10.50)''
  \begin{itemize}
  \item We already observed $\bar{x} = 8.9$; no probability involved
  \end{itemize}
\end{itemize}

\vspace{0.3cm}
\textbf{Bayesian methods} allow probability statements about parameters, but that's beyond this course.
\end{frame}

\begin{frame}
\frametitle{Other Confidence Levels}
\textbf{Not always 95\%!}

\vspace{0.3cm}
\begin{center}
\begin{tabular}{cccc}
\toprule
Confidence & $\alpha$ & Z critical & Interpretation \\
\midrule
90\% & 0.10 & 1.645 & Narrower, less confident \\
95\% & 0.05 & 1.96 & Standard choice \\
99\% & 0.01 & 2.576 & Wider, more confident \\
\bottomrule
\end{tabular}
\end{center}

\vspace{0.3cm}
\textbf{Trade-off:}
\begin{itemize}
\item Higher confidence $\Rightarrow$ wider interval (less precise)
\item Lower confidence $\Rightarrow$ narrower interval (more precise)
\item Choice depends on context and consequences of error
\end{itemize}

\vspace{0.3cm}
\textbf{General formula:}
$$\text{CI} = \bar{x} \pm z_{\alpha/2} \times \frac{\sigma}{\sqrt{n}}$$

where $z_{\alpha/2}$ leaves area $\alpha/2$ in each tail
\end{frame}

\section{R Examples and Practice}

\begin{frame}[fragile]
\frametitle{R Example: Binomial Distribution}
\small
\begin{verbatim}
# Vaccine efficacy: n=20, p=0.90
n <- 20
p <- 0.90

# Probability exactly 18 protected
dbinom(18, n, p)  # 0.2852

# Probability at least 18 protected
pbinom(17, n, p, lower.tail=FALSE)  # 0.6769
# or: 1 - pbinom(17, n, p)

# Expected value and variance
expected <- n * p           # 18
variance <- n * p * (1-p)   # 1.8
sd <- sqrt(variance)        # 1.34

# Visualize
x <- 0:n
probs <- dbinom(x, n, p)
barplot(probs, names.arg=x,
        main="Binomial(20, 0.90)",
        xlab="Number Protected", ylab="Probability",
        col="lightblue")
abline(v=expected+0.5, col="red", lty=2, lwd=2)
\end{verbatim}
\end{frame}

\begin{frame}[fragile]
\frametitle{R Example: Normal Distribution}
\small
\begin{verbatim}
# Cholesterol: mean=200, SD=40
mu <- 200
sigma <- 40

# Probability cholesterol > 240
pnorm(240, mu, sigma, lower.tail=FALSE)  # 0.159

# Probability between 180 and 220
pnorm(220, mu, sigma) - pnorm(180, mu, sigma)  # 0.383

# 90th percentile
qnorm(0.90, mu, sigma)  # 251.3 mg/dL

# Visualize with shaded regions
x <- seq(100, 300, length=200)
y <- dnorm(x, mu, sigma)
plot(x, y, type="l", lwd=2,
     main="Adult Cholesterol Distribution",
     xlab="Cholesterol (mg/dL)", ylab="Density")
abline(v=c(200, 240), col=c("blue", "red"), lty=2)
legend("topright", c("Mean", "High (>=240)"),
       col=c("blue", "red"), lty=2)
\end{verbatim}
\end{frame}

\begin{frame}[fragile]
\frametitle{R Example: Sampling Distribution Simulation}
\footnotesize
\begin{verbatim}
# Simulate CLT
set.seed(524)
pop_mean <- 120    # Population mean BP
pop_sd <- 20       # Population SD

# Different sample sizes
sample_sizes <- c(5, 10, 30, 100)

par(mfrow=c(2,2))
for(n in sample_sizes) {
  # Draw 1000 samples, compute means
  sample_means <- replicate(1000, {
    sample <- rnorm(n, pop_mean, pop_sd)
    mean(sample)
  })

  # Plot sampling distribution
  hist(sample_means, breaks=30, prob=TRUE,
       main=paste("n =", n),
       xlab="Sample Mean", col="lightblue")

  # Overlay theoretical normal
  se <- pop_sd/sqrt(n)
  curve(dnorm(x, pop_mean, se), add=TRUE, col="red", lwd=2)
}
\end{verbatim}
\end{frame}

\begin{frame}[fragile]
\frametitle{R Example: Confidence Intervals}
\small
\begin{verbatim}
# Tumor growth data (control group)
control <- c(7, 10, 9, 8, 7, 6, 8, 9, 12, 13)

# Method 1: Manual calculation
xbar <- mean(control)                    # 8.9
s <- sd(control)                         # 2.234
n <- length(control)                     # 10
se <- s/sqrt(n)                          # 0.707
t_crit <- qt(0.975, df=n-1)             # 2.262
margin <- t_crit * se                    # 1.598

ci_lower <- xbar - margin                # 7.30
ci_upper <- xbar + margin                # 10.50

cat("95% CI:", round(ci_lower, 2), "to",
    round(ci_upper, 2), "\n")

# Method 2: Using t.test() function
result <- t.test(control, conf.level=0.95)
result$conf.int  # (7.30, 10.50)
\end{verbatim}
\end{frame}

\begin{frame}[fragile]
\frametitle{R Example: Comparing Z and t CIs}
\footnotesize
\begin{verbatim}
# Sample data
data <- rnorm(15, mean=100, sd=15)
xbar <- mean(data)
s <- sd(data)
n <- length(data)
se <- s/sqrt(n)

# Z-based CI (assumes sigma known = 15)
z_crit <- qnorm(0.975)
z_ci <- xbar + c(-1, 1) * z_crit * 15/sqrt(n)

# t-based CI (sigma unknown, use s)
t_crit <- qt(0.975, df=n-1)
t_ci <- xbar + c(-1, 1) * t_crit * se

# Compare
cat("Z-based CI:", round(z_ci, 2), "\n")
cat("t-based CI:", round(t_ci, 2), "\n")
cat("t CI is", round((t_ci[2]-t_ci[1])/(z_ci[2]-z_ci[1]), 3),
    "times wider\n")

# Note: Difference decreases as n increases
# For n=100, they're nearly identical
\end{verbatim}
\end{frame}

\begin{frame}
\frametitle{Summary}
\textbf{Today we covered:}
\begin{itemize}
\item \textbf{Chapter 7.1:} Binomial distribution
  \begin{itemize}
  \item Conditions, probability formula, mean and variance
  \item R functions: dbinom(), pbinom()
  \end{itemize}
\item \textbf{Chapter 7.2:} Normal distribution
  \begin{itemize}
  \item Properties, standardization, empirical rule
  \item R functions: dnorm(), pnorm(), qnorm()
  \end{itemize}
\item \textbf{Chapter 7.4:} Other distributions (t, chi-square, F)
\item \textbf{Chapter 8:} Sampling distributions
  \begin{itemize}
  \item Central Limit Theorem
  \item Standard error vs standard deviation
  \item Confidence intervals (introduction)
  \end{itemize}
\end{itemize}

\vspace{0.5cm}
\textbf{Next Week:} Chapter 9 - More on Confidence Intervals

\vspace{0.3cm}
\textbf{Reading:} Complete all of Chapter 9
\end{frame}

\begin{frame}
\frametitle{Questions?}
\begin{center}
\Large{Thank you!}

\vspace{1cm}

\textbf{Email:} john.molitor@oregonstate.edu\\
\textbf{Office:} 153 Milam
\end{center}
\end{frame}

\end{document}
